% Canonical source: floquet_analysis/main.tex (inherited May 2025 draft; line locators from
% doc/auto/source-inventory.md and direct inspection: 66–140, 145–220, 295–345, 437–565, 675–750);
% doc/auto/research-program.md (corrections); doc/auto/graph-reproduction.md (2NH state, unverified multipliers).
% Full assertion-by-assertion table: doc/auto/theory-propositions.md §3 (this appendix defers to it).
\section{The inherited May 2025 draft and what replaces it}
\label{app:inherited}

The project inherited an unpublished note, \texttt{floquet\_analysis/main.tex} (May 2025), that analysed a GC-LSTM trained on the T-GRAB periodic task with Floquet terminology. The note is preserved unchanged as a source snapshot; its figures have no recovered checkpoint, raw arrays or producing code. \Cref{tab:inherited} lists what it asserted and which statement of this paper replaces each inference. The full assertion-by-assertion table lives in \texttt{doc/auto/theory-propositions.md}. \evidence{attributed-unverified}

\begin{table}[h]
\centering
\caption{Inherited draft assertions (with line locators in \texttt{floquet\_analysis/main.tex}) and their replacements.}
\label{tab:inherited}
\small
\begin{tabular}{p{0.07\linewidth}p{0.41\linewidth}p{0.44\linewidth}}
\toprule
Lines & Draft assertion & Replacement \\
\midrule
140 & Perfect accuracy on the periodic task implies the model internally encodes a modulo-$\tau$ specification, to be probed as Floquet modes. & Next-graph success does not establish retention of prefix information; node-identity features permit phase recognition from the current input. Required memory is tested with independently written cues (\cref{sec:design}). \\
160--210 & The GC-LSTM update is the Euler step of an ODE $\dot s=F(s,A)-s$, giving a ``continuous-time limit''. & The exact flow of that ODE is not $F$; no embedding is needed. All analysis is of the discrete map (\cref{sec:setting}). \\
216--219 & Periodic input implies a stable periodic reference solution $s^*(t+T)=s^*(t)$. & Closure must be verified at declared tolerances (\cref{def:verified-orbit}); affine forced systems do close (P5), nonlinear ones need not; the response period may be a multiple of the input period or absent. \\
299--340 & The monodromy matrix is the ordered product of per-step Jacobians. & Retained, conditional on a verified orbit; otherwise the product is a finite-time propagator (\cref{sec:setting}). \\
353, 484 & State is $(h,c)\in\R^{2H}$ and $J_k\in\R^{2H\times2H}$. & The complete graph state is $2NH$ with node coupling through graph convolution; a per-node block is not the transition Jacobian. \\
437--452 & With constant block Jacobians $J_A,J_B$, AABB yields $M=J_B^2J_A^2$ and predicts four undamped multipliers $e^{i2\pi m/4}$. & Repeated input does not imply equal Jacobians (states differ between repeats). A one-step rotation at frequency $m/T$ raised to $T$ has one-period multiplier $1$, so AABB predicts no fourth roots of unity for the one-period map. \\
502--510 & Six leading eigenvalues with $|\lambda_j|\approx0.41$--$0.44$ decay by $\approx0.42$ per cycle. & Unverified attributed observation (no checkpoint, arrays or code); a top-$K$ partial spectrum; even if reproduced it shows only local per-cycle contraction of those modes. \\
512--531 & Monotone envelope decrease over 20 steps confirms decay from $|\lambda|<1$. & $|\mu|<1$ does not imply monotone within-cycle decay and monotone decay does not identify the multiplier (P8). \\
541--565 & Small linearization error means the network ``converges to linear dynamics''. & Local first-order accuracy at a point is not global linearity; it validates the Jacobian, nothing more. \\
679--716 & A PSD peak at $f=1/4$ or $1/2$ would be a unit-modulus Floquet mode ``in disguise''; its absence shows no undamped clock, only re-encoding. & A periodic response with multiplier $a^T<1$ has power at the input frequencies (P5); a trajectory spectrum is not a tangent spectrum in either direction. \\
730--747 & A zig-zag PCA trajectory indicates discrete re-encoding without a smooth orbit. & A three-dimensional projection does not identify an internal mechanism or exclude a clock. \\
\bottomrule
\end{tabular}
\end{table}

The two inherited benchmark manuscripts (\texttt{ACM\_\_\_TGRAB/}, \texttt{ACM\_\_\_TGRAB\_KDD\_2026/}) are likewise preserved; the entrypoint \texttt{ACM\_\_\_TGRAB\_KDD\_2026/kdd/main.tex} (lines 327--354) selects the \texttt{\_ai} section variants and the appendix. The task notation differs between the drafts ($P_{\mathrm{det}}(k,n)$ with $k$ patterns and block length $n$ in the benchmark; $N$ patterns and block length $k$ in the Floquet note); AABB is $P_{\mathrm{det}}(2,2)$ and AB is $P_{\mathrm{det}}(2,1)$, and which of them the inherited figures used is unresolved. \todo{Insert the resolved task identity and recovered provenance, or record that recovery failed and label the graph work a reconstruction.}
